\section{From the mean arithmetic gap to the critical line}
\label{sec:conditional-main-v42}

Fix a primitive Dirichlet character $\psi$. To rule out a zero to the
right of the critical line we shall bound a weighted sum of primes on
intervals $[Y,Y^{3/2}]$ with $Y$ arbitrarily large. The arithmetic
construction produces such an interval and, for every real $x$ in it,
a collection of theta periods. Two properties of these periods matter
here. Their finite data pick out the primitive characters inducing the
powers of $\psi$. Their integral and complex realizations together give
a bound on a mean over translates. We state both properties precisely
before deducing the bound for primes.

% excerpt: \subsection{The fixed analytic target} || % END CMP6

\subsection{How the periods bound the prime sum}

The prime-power sums $\mathcal A(a_t,t,w_{u,j})$ in (CMP.8) record the
divisibility forced by the derivatives of integral theta sections.
The problem is to bound how far their sum can exceed the sum of the
explicit main terms $\operatorname{Re}\mathcal M(a_t,w_{u,j})$. To see
where the complex periods enter, take a nonzero minor $\Delta$ of one
of the period matrices, expressed in integral bases over its coefficient
field $K$. Factoring the norm of $\Delta$ gives
\[
 \sum_{\mathfrak p} v_{\mathfrak p}(\Delta)\log N\mathfrak p
   =\sum_{\sigma:K\hookrightarrow\mathbb C}\log|\sigma(\Delta)|.
\]
A lower bound on the left can therefore be compared with an upper bound
on the right, provided the complex calculation uses the images of the
same rows in the same bases. In the slope calculation the norms of the
exterior powers keep track of the changes of metric.

There are two losses to control before applying this argument. Clearing
denominators costs part of the divisibility obtained at finite primes.
Also, exceptional rows have to be discarded. The surviving images must
span enough independent directions for the slope inequality to apply.
The diagonal and mixed second moments of the periods bound this loss of
rank. With the rank established, the height and norm bounds on the
surviving maps enter the slope calculation. The
\hyperref[sec:tied-window-completion]{calculation for one translate}
gives the error budget $x^{C\varepsilon}B$ with $B$ as in (CMP.7).
The same bound holds on taking the mean over $u$. This gives precisely
the estimate for $\overline\Gamma$ used in (CMP.11).

The choice of interval is part of the arithmetic argument. The auxiliary
primes and the field must allow the congruence estimates that control
the denominators, as well as the character estimates in (CMP.4). For
general $\psi$ one must realize the primitive inducers of its powers
on this field. The construction uses a chosen residual representation
and propagates it through a chain of rank-four representations. For
$\psi=1$ the signal in (CMP.1) is already $G$. In both cases we fix
$\psi$ and $\varepsilon$ before choosing $Y$ so that the constants in
the estimates work on arbitrarily late intervals.

The \hyperref[ms-manuscript-notes-arithmetic-common-design-window-transfer-md]{common-window
construction} gives the order of the geometric choices. The
\hyperref[ms-manuscript-notes-theta-centered-period-rank-and-slope-md]{rank and slope argument}
shows how the period estimates are combined. For the underlying formal
declarations and the earlier proofs of the analytic steps, see the
\href{https://github.com/bmatschke/grh-alpoge-claude/blob/edition-5/research/modules/construction/conditional_main_v42/contracts.json}{source concordance}.

\section{Vanishing of the polarized adjoint Selmer group}
\label{sec:adk-match}

One cohomological input to the congruence argument can be obtained
directly from Thorne's vanishing theorem. The application uses the
full polarized adjoint representation, including its scalar line. We
identify the coefficient fields, the Galois actions and the local
conditions before applying the theorem. The rank-three case is the one
needed for the Kato argument. The comparison below holds for all
$n\geq2$, and its last part explains the further local condition
required for rigidity on the chosen host.

% excerpt: \subsection{The exact statement to be supplied} || EOF
